\SourceLeaf{532}
\fancyhead[L]{\small CONNAISSANCE DES TEMS.}
\begin{center}\small Göttingische gelehrte Anzeigen. Stück 153. S. 1526..1528. 24. September 1810.\\[3pt]\rule{.72\linewidth}{.3pt}\end{center}

\emph{Connaissance des tems ou des mouvemens célestes, à l'usage des astronomes et des navigateurs, pour l'an} 1811, \emph{publiée par le bureau des longitudes.} 1809. Octav. Paris.

Wir übergehen den Kalender, dessen Einrichtung unverändert geblieben ist, und zeigen nur den Inhalt der Zusätze an. Drei Viertel davon bestehen aus den astronomischen Beobachtungen, die 1807 und 1808 an der kaiserlichen Sternwarte in Paris angestellt wurden. Ihre Bedeutung brauchen wir hier nicht eigens zu rühmen. Außer der Sonne, dem Mond und den wichtigsten Fixsternen wurden alle älteren Planeten sehr fleißig beobachtet. Von den neuen finden sich für 1807 ebenfalls zahlreiche Beobachtungen, besonders von Vesta und Ceres. Pallas wurde mehrmals beobachtet, Juno gar nicht. Für 1808 finden wir unter den neuen Planeten nur Vesta und Juno beobachtet, Juno nur einmal. In diesem Beobachtungsjournal fallen uns die mitunter sehr beträchtlichen Unregelmäßigkeiten im Gang der Pendeluhr auf. Man suchte ihnen durch häufigeres Verschieben der Pendellinse abzuhelfen. Mit dem parallaktischen Instrument wurden Vesta und der große Komet von 1807 beobachtet. Für den Kometen teilt \textsc{Bouvard} seine Bahnelemente mit. Darauf folgt ein Aufsatz von \textsc{Laplace} über die Abnahme der Schiefe der Ekliptik nach alten Beobachtungen. Besonders interessant ist hier die Erörterung einiger alter chinesischer Beobachtungen. Die älteste davon wurde 1100 Jahre vor unserer Zeitrechnung gemacht. Alle stimmen, ebenso wie die griechischen, arabischen und persischen Bestimmungen, mit der von \textsc{Laplace} aus der Theorie ermittelten Abnahme so gut überein, wie man es nur erwarten kann. Ein anderer Aufsatz von \textsc{Laplace} behandelt die Rotation des Saturnrings. Dazu haben bekanntlich \textsc{Schröter}s Beobachtungen ein von denen \textsc{Herschel}s völlig verschiedenes und mit der Theorie unvereinbares Ergebnis geliefert. Nach den ersteren schienen auf dem Ring Unregelmäßigkeiten vorhanden zu sein, die ihre Stelle nicht änderten. \textsc{Laplace} erklärt diese Beobachtungen durch die einfallsreiche Hypothese, dass der Ring eigentlich aus vielen konzentrischen Ringen besteht, die in verschiedenen Ebenen liegen und durchaus solche Unregelmäßigkeiten zeigen konnten. Solange sich die Lage der Ebenen der einzelnen Ringe nicht merklich änderte, würden diese Unregelmäßigkeiten natürlich unverändert erscheinen, auch wenn die Ringe rotierten. Unerklärt bleibt dabei nur, warum die beiden Ansen nicht denselben Anblick boten. Man darf hoffen, dass eine künftige Wiederholung dieser heiklen Beobachtungen darüber Aufschluss geben wird. — Eine darauf folgende Abhandlung von \textsc{Delambre} über die Bestimmung der Breite und der Zeit aus den beobachteten Höhen zweier bekannter Sterne ist hauptsächlich als Kommentar zu einem auch in unseren Blättern angezeigten Programm von \textsc{Gauss} gedacht (s. Gött. gel. Anz. 1808, S. 1945 [S. 37 dieses Bandes]). Allerdings scheint \textsc{Delambre} diese kleine Schrift nur aus dem Auszug in der \emph{Monatlichen Correspondenz} gekannt zu haben und beurteilt deshalb ihren Zweck nicht ganz richtig. \textsc{Delambre} führt auch die von \textsc{Mollweide} in der \emph{Monatl. Correspondenz} mitgeteilte Lösung desselben Problems an. Obwohl er sie aus anderen Gründen hinter die trigonometrische Lösung stellt, rühmt er die Leichtigkeit, mit der sich bei ihr aus den beiden Wurzeln der Gleichung die richtige auswählen lasse. Diesem Urteil können wir nicht beipflichten. Das von \textsc{Mollweide} angeführte Kriterium ist zwar einfach, aber unzulässig und der Aufgabe fremd. — Berechnung der Beobachtungen der Sonnenfinsternis vom 16. Juni 1806 in Utrecht, Mailand, München und Lilienthal von \textsc{van Beek Calkoen}. — \textsc{Burckhardt} über die verschiedenen Arten, wie Astronomen die Sonne beobachtet haben. — Von demselben Astronomen noch drei kleine Aufsätze über die Kometen von 1701 und 1772 und über das \SourceLeaf{533}Rostpendel. — Von Werken aus dem Gebiet der Astronomie ist diesmal nur die Reise von \textsc{d’Entrecasteaux} angezeigt. — Ein Auszug aus dem meteorologischen Journal, das \textsc{Bouvard} 1807 auf der Sternwarte führte, bildet den Abschluss.
